\section{Conclusion}
\label{sec:conclusion}

\system{} tested a simple but frequently skipped premise: before routing among RAG
and memory actions, do those actions produce generator-specific, support-qualified
quality separation at a measurable cost? In the frozen benchmark, the answer was
mixed. An identical grounded route materially improved Granite-4.1-3B and Qwen3-4B,
had an uncertain positive effect for SmolLM3, removed the only successful direct
stratum for Qwen3-0.6B, and exposed a deterministic integration failure for
Granite-4.0-1B. Every grounded endpoint had a higher observed generation-window
GPU-board energy mean, with four matched pairs measured on the same board.

The router consequently had no validation-feasible operating point. Its disclosed
fail-closed policy used the same fallback for every test query, achieved zero strict
success, and reduced energy by exactly zero; the confirmatory hypothesis failed.
This negative result is more useful than an optimistic aggregate. It identifies
route viability, interface compliance, task-conditioned grounding value, precise
measurement boundaries, and nonzero clean utility as prerequisites for credible
energy-aware RAG routing.

Future work should repair model-specific interfaces under a new frozen specification,
add public and long-context benchmarks, measure whole-system energy with repeated
interleaved runs, and train a router only after its calibration split demonstrates
a feasible operating point and query-dependent separation among compatible routes.
Until then, retrieval and memory should be treated as generator-dependent
interventions whose value must be measured, not assumed.
